% ===============================
% Worksheet / Manual Template
% ===============================
\documentclass[12pt, a4paper]{article}

% ===============================
% Metadata
% ===============================
\newcommand*{\CourseCode}{MAT321}
\newcommand*{\CourseTitle}{Real Analysis II}
\newcommand*{\Chapter}{Compactness in Metric Spaces}
\newcommand*{\Topics}{Heine-Borel Theorem and Extreme Value Theorem}
\newcommand*{\DocDate}{April 24, 2026}

\include{header}

% ==========================================
% Document
% ==========================================
\begin{document}
\FrontPage
\Large

\begin{heading}
    Open Cover and Subcover
\end{heading}

\vspace{10pt}

\begin{definition}[Open Cover]
    Let $A$ be a subset of a metric space $(X, d)$. An \textit{open cover} of $A$ is a collection $\{U_\alpha\}_{\alpha \in I}$ of open subsets of $X$ such that
    \[
        A \subseteq \bigcup_{\alpha \in I} U_\alpha.
    \]
\end{definition}

\clearpage

\begin{definition}[Subcover]
    A \textit{subcover} of an open cover $\{U_\alpha\}_{\alpha \in I}$ of a set $A$ is a subcollection $\{U_\alpha\}_{\alpha \in J}$, where $J \subseteq I$, such that
    \[
        A \subseteq \bigcup_{\alpha \in J} U_\alpha.
    \]
\end{definition}

\clearpage

\begin{definition}[Finite Subcover]
    A \textit{finite subcover} of an open cover $\{U_\alpha\}_{\alpha \in I}$ of a set $A$ is a subcollection $\{U_\alpha\}_{\alpha \in J}$, where $J \subseteq I$ is a finite set, such that
    \[
        A \subseteq \bigcup_{\alpha \in J} U_\alpha.
    \]
\end{definition}

\clearpage

\begin{heading}
    Compactness
\end{heading}

\begin{definition}[Compactness]
    A subset $A$ of a metric space $(X, d)$ is called \textit{compact} if every open cover of $A$ has a finite subcover. That is, if $\{U_\alpha\}_{\alpha \in I}$ is an open cover of $A$, then there exists a finite subset $J \subseteq I$ such that
    \[
        A \subseteq \bigcup_{\alpha \in J} U_\alpha.
    \]
\end{definition}

\clearpage

\begin{problem}
    Show that the Real line $\mathbb{R}$ is not compact in the standard metric.
\end{problem}

\clearpage

\begin{problem}
    Show that the closed interval $(0, 1)$ is not compact in the standard metric on $\mathbb{R}$.
\end{problem}

\clearpage

\begin{problem}
    Show that any finite subset of real numbers is compact in the standard metric on $\mathbb{R}$.
\end{problem}

\clearpage

\begin{problem}
    Show that the closed interval $[0, 1]$ is compact in the standard metric on $\mathbb{R}$.
\end{problem}

\clearpage

\begin{heading}
    Why Compactness is Important?
\end{heading}

\clearpage

\begin{heading}
    Boundness of Functions
\end{heading}

\vspace{10pt}

\begin{theorem}[Boundedness of Continuous Functions on Compact Sets]
    Let $(X, d)$ be a metric space and let $f : X \to \mathbb{R}$ be a continuous function. If $K \subseteq X$ is compact, then $f$ is bounded.
\end{theorem}

\clearpage

\begin{heading}
    Some Important Theorems on Compactness
\end{heading}

\vspace{10pt}

\begin{theorem}
    Every closed subset of a compact metric space is compact.
\end{theorem}

\clearpage

\begin{theorem}
    Every compact subset of a metric space $(X, d)$ is closed.
\end{theorem}

\clearpage

\begin{theorem}
    Every compact subset of a metric space $(X, d)$ is bounded.
\end{theorem}

\clearpage

\begin{heading}
    Heine-Borel Theorem
\end{heading}

\vspace{10pt}

\begin{theorem}[Heine-Borel Theorem]
    A subset of $\mathbb{R}$ is compact if and only if it is closed and bounded.
\end{theorem}

\clearpage

\begin{heading}
    Continuous Image and Extreme Value Theorem
\end{heading}

\begin{theorem}[Continuous Image of Compact Set is Compact]
    Let $(X, d)$ be a metric space and let $f : X \to \mathbb{R}$ be a continuous function. If $K \subseteq X$ is compact, then $f(K)$ is compact in $\mathbb{R}$.
\end{theorem}

\clearpage

\begin{heading}
    Extreme Value Theorem
\end{heading}

\vspace{10pt}

\begin{theorem}[Extreme Value Theorem]
    Let $(X, d)$ be a metric space and let $f : X \to \mathbb{R}$ be a continuous function. If $K \subseteq X$ is compact, then $f$ attains its maximum and minimum values on $K$.
\end{theorem}


\EndPage
\end{document}